\section{Introduction}

Public organizations sometimes take rare discretionary actions in response to incidents, releases,
or salient milestones. Such actions invite retrospective explanations, but a useful scientific
question is narrower: can information that was publicly available at time $t$ improve a calibrated
forecast of an action after $t$? Answering this question requires more than a classifier. The event
definition, source hierarchy, observation boundary, missing runs, and later data corrections must
all be auditable.

We examine public announcements of special Codex usage-limit resets. The setting is deliberately
small and difficult. Events are irregular, the policy regime changes, multiple announcements can
refer to one operational action, and reasons are often revealed in the announcement itself. Treating
those reasons as predictors would create direct leakage, a general threat to machine-learning
science \citep{kapoor2023leakage}. The project therefore forecasts public organizational behavior
rather than inferring a person's private motives or personality.

Our contributions are fourfold. First, we define an announcement-level gold standard and convert it
into daily and six-hour person-period datasets with explicit $(\text{start},\text{end}]$ boundaries.
Second, we compare global, adaptive-rate, renewal, calendar, and official-incident models under a
common expanding-window evaluation. A rolling 30-day rate defeats the initially selected calendar
model, illustrating why weak global-rate baselines are insufficient. Third, official-incident
features do not improve reliably over temporal baselines. Fourth, we implement a prospective protocol
with hash-locked predictions, independent outcome tables, explicit exclusions, missed-run records,
and append-only revisions.

The historical experiments are model-development evidence rather than a prospective test. The
confirmatory sequence begins only with valid forecasts issued at the frozen daily time. This
distinction prevents a favorable historical result from being presented as real-time performance.
